\begin{textsample}{NEG Dim 2 – persona\_grok – Score -1.00 – t173\_grok.txt}  \label{ex:f2_neg_025}
Fast \textbf{fashion} often touts itself as circular, but this is a myth.
Overproduced \textbf{clothes} from the Global North are exported as waste to the Global South, where they pollute environments and harm health.
In Kenya, imports of these \textbf{clothes} reached 185,000 tonnes in 2019, with 30-40 % of them unusable.
This results in 150-200 tonnes being dumped or burnt every day, as observed in places like Dandora and along the Nairobi River.
Vendors who deal with bales of Mitumba frequently find them filled with poor-quality waste.
To address this issue, \textbf{brands} need to produce fewer \textbf{clothes} that are durable and repairable.
They must also end neocolonial dumping practices and support local manufacturing.
While the EU strategy offers some help, a global treaty is essential to make real progress.

% matched lemmas: brand, clothes, fashion
\end{textsample}
